% page 526 du fac-simile (image p-525 du scan)
% bloc 160.0 x 263.6 mm ; marge gauche 8.0 mm
\pgc{-12.700mm}{0.000mm}{
\l{}
\l{}
\l{}
\l{}
\l{\cel{3}518.}
\l{}
\l{\sou{senato}. I. (en\cel{1}\sou{Roma, antique}). Unesma (hierarkiale) korpo}
\l{\cel{3}politikala, ek la patrici. - II. (\sou{nun}). Unesma (hierar-}
\l{\cel{3}kiale) korpo politikala di ula stati moderna. - DEFIRS.}
\l{}
\l{\sou{senatuskonsulto}. I. (en\cel{1}\sou{Roma, antique)}. Rezolvo agita da la}
\l{\cel{3}senato. - II. (en\cel{1}\sou{Francia}).Rezolvo agita da la senato di}
\l{\cel{3}Francia lor la konsulo-rejimo e lor la imperio di Napole-}
\l{\cel{3}ono Iesma. - EF.}
\l{}
\l{\sou{senco}. Maniero ye qua kozo esas komprenenda. La\cel{1}\sou{senco}\cel{1}konsis-}
\l{\cel{3}tas ye la idei diversa quin lu expresas od implikas- kontre}
\l{\cel{3}ke lua\cel{1}\sou{signifiko}\cel{1}konsistas ye la objekti diversa a qui lu}
\l{\cel{3}esas aplikebla.-EFIS.}
\l{}
\l{\sou{sendar}. (\sou{trans., ad)}.\cel{2}Igar (ulo, ulu) departar a destinario.}
\l{\cel{3}- ED.}
\l{}
\l{\sou{senecio}. (\sou{bot.}) Planto ek la familio \textquotedbl{}kompozaji), di qua la}
\l{\cel{3}semino uzesas por nutrar la uceli. - L. senecio.- FL.}
\l{}
\l{\sou{seneshalo}. (\sou{lor la feudismo-rejimo, en Francia}).Oficiro feuda-}
\l{\cel{3}la qua chefdirektis la domego, komandis la armeo, judiciis,}
\l{\cel{3}e c. - DEFIRS.}
\l{}
\l{\sou{senila}. (\sou{medic.}) Di qua la korpo e la spirito divenis febla}
\l{\cel{3}pro oldeso.--DEFIS.}
\l{}
\l{\sou{seniora}. I. Qua naskis ante sua frato o frati. - II.\cel{1}\sou{senioro}:}
\l{\cel{3}persono plu evoza kam altra. - DEF.}
\l{}
\l{\sou{senso.}\cel{1}Fakultato (che la homo, che la animalo) perceptar sua}
\l{\cel{3}impresesi de la objekti materiala (tusho, vido, audo, gusto,}
\l{\cel{3}flaro). - DEFIRS.}
\l{}
\l{\sou{sensaciono}. Intereseso o sentiment-eciteso tre forta, che tur-}
\l{\cel{3}bo. - DEFIS.}
\l{}
\l{\sou{sensitivo}. (\sou{bot.)}\cel{2}Varietato di mimozo, di qua la folii su}
\l{\cel{3}retrofaldas kande onu tushas li. - L.\cel{1}\sou{mimosa pudica}. EF.}
\l{}
\l{\sou{sensu}ala.\cel{2}Qua ambicias o diletas la plezuro de la sensi.}
\l{\cel{3}-EFIS.}
\l{}
\l{\sou{sentar}. (\sou{trans.}) Impresesar agreable o dolorigante (a) per}
\l{\cel{3}sua sensi, (b) sen sua sensi. - EFIS.}
\l{}
\l{\sou{sentenco}. Opiniono, expresata per formulo dogmatala. -DEFIRS.}
\l{}
\l{\sou{sentimento}. Psiko-movo, psiko-reakto, en qua ne implikesas}
\l{\cel{3}la sensi. - EFIRS.}
\l{}
\l{\sou{sentinelo}. I. Soldato qua havas kom tasko surveyar avan pos-}
\l{\cel{3}teno armeala (kazerno, fortreso, fortifikajo, kampeyo)}
\l{\cel{3}por ke olca ne surprizesez.- II. (\sou{metaf.) Sentinelo:}La}
\l{\cel{3}persono qua surveyas por evitar surprizo.- EFIS.}
}
